\documentclass[10pt,a4paper]{amsart}
\usepackage[margin=19mm]{geometry}
\usepackage{amsmath,amssymb,mathtools}
\usepackage[scr=rsfs,cal=euler]{mathalfa}
\usepackage{xfrac}
\usepackage[unicode,hidelinks,bookmarks=false]{hyperref}
\newcommand{\ie}{i.e.\ }
\newcommand{\cf}{cf.\ }
\newcommand{\resp}{respectively\ }
\newcommand{\Subsection}{\subsection}
\providecommand{\nobreakdash}{\nobreak\hbox{-}}
\setlength{\emergencystretch}{2em}
\multlinegap0pt
\usepackage{xcolor}
\usepackage{float}
\usepackage{amssymb}
\usepackage{tikz}
\newtheorem{lemma}{Lemma}
\newcommand{\field}{K}
\renewcommand{\hom}{\mathrm{Hom}}
\title{DEGENERATION ORDER OF $3\times 3$ NILPOTENT MATRIX TUPLES}
\author{M\'aty\'as Domokos, Botond Mikl\'osi}
\date{Source: arXiv:2604.22609v1 (CC BY 4.0)}
\begin{document}
\maketitle

\noindent\emph{Source and licence.} DEGENERATION ORDER OF $3\times 3$ NILPOTENT MATRIX TUPLES. Version arXiv:2604.22609v1 (CC BY 4.0), distributed by arXiv under the Creative Commons Attribution 4.0 International licence, http://creativecommons.org/licenses/by/4.0/, as recorded in the arXiv licence metadata for this exact version and shown on the abstract page https://arxiv.org/abs/2604.22609v1. Full licence notice: \texttt{licenses/arxiv-2604.22609-CC-BY-4.0.txt}.\par\smallskip

\noindent\textit{Editorial scope.} Counted unit: the article's lemma lemma:S,T,M (source0.tex lines 450--461). The article's complete original proof of it is delivered with it (source0.tex lines 463--519, one \texttt{proof} environment of 57 lines). No mathematics is written, repaired or re-derived: the statement and the proof are the article's own lines, in the article's own notation, and no retained line is substituted or re-flowed.  No further source-local context is needed: every cross-reference of every retained line resolves inside the two delivered spans.  Nothing was omitted from any retained span, so the package carries no omission record.  No retained line contains a citation, so the wrapper carries no bibliography and no bibliography entry.  No retained line contains a figure environment, an image-inclusion command or drawing code, so no image is delivered and the package declares no asset dependency.  Editorial formatting only, in addition: an AMS \texttt{amsart} preamble that restates the article's own declarations and macros used by the retained lines, namely the environments \texttt{lemma} on separate counters, together with the standard packages those lines need.  The retained environments print in this document as Lemma 1 in order of appearance, whereas the article prints the same items with the numbering of its own full document.  The complete native source of the article as distributed in the arXiv e-print for this exact version is bundled unchanged as \texttt{sources/199110/source0.tex}.
\begin{lemma}\label{lemma:S,T,M} 
    Let $M$ and $N$ be $R$-modules, $\dim_\field M=\dim_\field N$.  
    Assume that $M$ 
    or $N$ has 
    Jordan-Hölder composition length at most $2$, and 
    \begin{equation}\label{eq:hom(X,M)}
      \text{for all submodules }X \text{ of }M\colon \quad   \dim_{\field}\hom_{R}(X, M) \leq \dim_{\field}\hom_{R}(X, N).
    \end{equation}    
    Then $M\cong N$, or $M$ has Jordan-Hölder composition length $2$, and 
    denoting by $S,T$ the composition factors of $M$  ($S\cong T$ is allowed), we have 
    $N\cong S\oplus T$. 
\end{lemma}

\begin{proof}
By Schur's Lemma we have that for a finite dimensional simple $R$-module $X$ and 
for a finite dimensional $R$-module $V$, $\dim_{\field}\hom_{R}(X, V)$ equals 
the multiplicity of $X$ as a direct summand in the socle of $V$. 
Hence applying \eqref{eq:hom(X,M)} for the simple submodules $X$ of $M$, 
we get that $N$ contains a submodule isomorphic to $\mathrm{Soc}(M)$. 
In particular, if $M=\mathrm{Soc}(M)$ is semisimple, then $\mathrm{Soc}(N)$ 
contains a direct summand isomorphic to $M$, implying in turn 
by $\dim_{\field}M=\dim_{\field}N$  that $M\cong N$. 
From now on we assume that $M$ is not semisimple, 
and $M\ncong N$. 

Suppose that $M$ has Jordan-Hölder composition length $2$. 
Then $\mathrm{Soc}(M)\cong S$ and 
$M/\mathrm{Soc}(M)\cong T$ are simple $R$-modules. 
As we explained above, $N$ has a submodule isomorphic to $S$. 
As $\dim_{\field}\hom_R(M,M)>0$, there exists a non-zero 
$\varphi\in \hom_R(M,N)$ by \eqref{eq:hom(X,M)} applied with $X=M$. 
As $\varphi$ is not an isomorphism,  
it has a non-zero kernel, so we have $\ker(\varphi)\supseteq \mathrm{Soc}(M)$, 
and the image of $\varphi$ must be isomorphic to $T$. So $N$ also has a submodule 
isomorphic to $T$. 
If $S\ncong T$, then $N$ has a submodule isomorphic to $S\oplus T$, 
and taking into account 
$\dim_{\field}M=\dim_{\field}S+\dim_{\field}T=\dim_{\field}N$  
we conclude that $N\cong S\oplus T$. 
If $S\cong T$, then . 
$M/\mathrm{Soc}(M)\cong S$.  
Denoting by $\varphi$ 
the natural surjection $M\to M/\mathrm{Soc}(M)$ composed with the 
embedding $\mathrm{Soc}(M)\to M$, we have that $\varphi$ and 
$\mathrm{id}_M$ are linearly independent elements in 
$\hom_R(M,M)$. It follows that 
$2\le \dim_{\field}\hom_R(M,M)\le 
\dim_{\field}\hom_R(M,N)$. 
Since $M\ncong N$, any homomorphism from $M$ to $N$ contains 
$\mathrm{Soc}(M)$ in its kernel, therefore 
$\dim_{\field}\hom_R(S,N)=\dim_{\field}\hom_R(M/\mathrm{Soc}(M),N)=
\dim_{\field}\hom_R(M,N)\ge 2$. 
Thus the multiplicity of $S$ as a direct summand in $N$ is at least $2$, 
implying in turn that $N\cong S\oplus S$. 

Suppose finally that $N$ has Jordan-Hölder composition length at most $2$. 
Since $N$ contains a submodule isomorphic to 
$\mathrm{Soc}(M)$ but 
$N\ncong \mathrm{Soc}(M)$ 
(since otherwise $N\cong M$ by $\dim_\field N=\dim_\field M$), 
we conclude that $\mathrm{Soc}(M)=S$ is simple, and $M$ has a submodule $M'$ 
with $M'\supseteq \mathrm{Soc}(M)$ and $M'/\mathrm{Soc}(M)=T$ simple. 
Then $1\le \dim_\field\hom_R(M',M)\le \dim_\field\hom_R(M',N)$. 
So there exists a non-zero homomorphism $\varphi\in \hom_R(M',N)$. 
It can not be an isomorphism, hence it factors through $T$. 
Moreover, if $S\cong T$, then $\dim_\field\hom_R(M',M)\ge 2$. 
So in the same way as in the above paragraph, we conclude that $N\cong S\oplus T$, which implies 
in turn by $\dim_\field M=\dim_\field N$ 
that $M=M'$. 
\end{proof}

\end{document}
